% ============================================================
% CHAPTER 22: STEINER NETWORK (matches Vazirani Ch. 22/23)
% ============================================================
\setcounter{chapter}{21}
\chapter{Steiner Network}
\label{ch:steiner_network}

\section{Intuition: LP Extreme Points and Iterative Rounding}

\begin{intuitionbox}
\textbf{Peeling and Layering Intuition:} Primal-dual schemas and randomized rounding are powerful, but some problems require a more direct structural exploitation of the LP relaxation. Kamal Jain (1998) introduced \emph{Iterative Rounding} for the Survivable Network Design Problem (SNDP). The key insight is that any extreme point (basic feasible solution) of the SNDP LP relaxation must contain at least one edge $e$ with fractional value $x_e^* \ge 1/2$. By rounding this edge up to 1, reducing the connectivity requirements, and solving the LP on the remaining graph iteratively, we obtain a guaranteed 2-approximation.
\end{intuitionbox}

\section{Problem Definition}
Given an undirected graph $G = (V, E)$ with edge costs $c_e \ge 0$, and connectivity requirements $r_{uv} \ge 0$ for each pair of vertices $u, v \in V$, the \emph{Steiner Network} (or Survivable Network Design) problem is to find a minimum cost subgraph $H \subseteq E$ containing at least $r_{uv}$ edge-disjoint paths between $u$ and $v$ for all $u, v \in V$.

The problem is NP-hard. We seek a 2-approximation using Jain's iterative rounding algorithm.

\section{Algorithm and Python Implementation}
We solve the LP relaxation of SNDP by formulating its dual (which has a naturally feasible origin since edge costs are non-negative) and iteratively round edges with $x_e \ge 0.5$.

\begin{lstlisting}[style=python, caption=Jain's Iterative Rounding Algorithm]
class Simplex:
    """Maximize c x  s.t. Ax <= b, x >= 0."""
    def __init__(self, A, b, c):
        self.m = len(A)
        self.n = len(A[0]) if A else 0
        self.A = [row[:] for row in A]
        self.b = b[:]
        self.c = c[:]
        self.obj_row = []

    def solve(self):
        tab = [self.A[i] + [1 if j == i else 0 for j in range(self.m)]
               + [self.b[i]] for i in range(self.m)]
        tab.append([-v for v in self.c] + [0] * self.m + [0])
        while True:
            e = min(range(self.n + self.m), key=lambda j: tab[-1][j])
            if tab[-1][e] >= -1e-9:
                break
            l = min(range(self.m), key=lambda i: tab[i][-1] / tab[i][e]
                    if tab[i][e] > 1e-9 else float('inf'))
            if tab[l][e] <= 1e-9:
                return None, float('inf')
            p = tab[l][e]
            tab[l] = [v / p for v in tab[l]]
            for i in range(self.m + 1):
                if i != l:
                    f = tab[i][e]
                    tab[i] = [tab[i][j] - f * tab[l][j]
                              for j in range(len(tab[0]))]
        x = [0.0] * self.n
        for i in range(self.m):
            for j in range(self.n + self.m):
                if (abs(tab[i][j] - 1) < 1e-6 and
                        sum(abs(tab[r][j]) for r in range(self.m + 1)) < 1.1):
                    if j < self.n:
                        x[j] = max(0.0, tab[i][-1])
        self.obj_row = tab[-1]
        return x, tab[-1][-1]


def get_all_cuts(n):
    """Every non-empty proper vertex subset."""
    return [{j for j in range(n) if i & (1 << j)}
            for i in range(1, 1 << (n - 1))]


def solve_sndp_lp_phase(n, edges, costs, r, fixed_edges):
    """LP relaxation for SNDP via its dual (global cut constraints)."""
    cuts = get_all_cuts(n)
    n_cuts, n_edges = len(cuts), len(edges)
    active = [i for i in range(n_edges) if i not in fixed_edges]
    n_active = len(active)

    f = []
    for cut in cuts:
        val = max([req for (u, v), req in r.items()
                   if (u in cut) != (v in cut)] + [0])
        fixed_cross = sum(1 for idx in fixed_edges
                          if (edges[idx][0] in cut) != (edges[idx][1] in cut))
        f.append(max(0, val - fixed_cross))

    n_cols = n_cuts + n_active
    A_dual, b_dual = [], []
    for j, ei in enumerate(active):
        u, v = edges[ei]
        row = [0.0] * n_cols
        for s_idx, cut in enumerate(cuts):
            if (u in cut) != (v in cut):
                row[s_idx] = 1.0
        row[n_cuts + j] = -1.0
        A_dual.append(row)
        b_dual.append(costs[ei])

    solver = Simplex(A_dual, b_dual, [float(v) for v in f] + [-1.0] * n_active)
    dual_sol, _ = solver.solve()
    if dual_sol is None:
        return [0.0] * n_edges

    x = [0.0] * n_edges
    for i in fixed_edges:
        x[i] = 1.0
    for j, ei in enumerate(active):
        x[ei] = max(0.0, solver.obj_row[n_cols + j])
    return x


def jain_iterative_rounding(n, edges, costs, r):
    """Jain's iterative rounding: round an edge x_e >= 1/2 per LP phase."""
    fixed_edges = set()

    def requirements_met():
        for cut in get_all_cuts(n):
            req = max([val for (u, v), val in r.items()
                       if (u in cut) != (v in cut)] + [0])
            if req == 0:
                continue
            cap = sum(1 for i in fixed_edges
                      if (edges[i][0] in cut) != (edges[i][1] in cut))
            if cap < req:
                return False
        return True

    while not requirements_met():
        x = solve_sndp_lp_phase(n, edges, costs, r, fixed_edges)

        active = [i for i in range(len(edges)) if i not in fixed_edges]
        best_idx = max(active, key=lambda i: x[i])
        fixed_edges.add(best_idx)

    return [edges[i] for i in fixed_edges]
\end{lstlisting}

\section{Analysis: Jain's Theorem and Factor 2}
The LP relaxation for the Steiner Network problem is:
\[
\min \sum_{e \in E} c_e x_e \quad \text{s.t.} \quad \sum_{e \in \delta(S)} x_e \ge f(S) \; (\forall S \subset V), \quad 0 \le x_e \le 1
\]
where $f(S) = \max_{u \in S, v \notin S} r_{uv}$.
An extreme point (or basic feasible solution) of this LP is defined by a set of linearly independent tight cut constraints. Jain proved a fundamental structural property of these extreme points:
\begin{theorem}[Jain 1998]
For any connectivity requirement function $f$ that is weakly supermodular, any basic feasible solution to the LP relaxation has at least one edge $e$ with $x_e \ge 1/2$.
\end{theorem}
Based on this theorem, we round $x_e$ to 1, which increases its cost to $c_e \le 2 \cdot c_e x_e$. By induction on the iterations, the total cost of the rounded edges is bounded by:
\[
\sum_{e \in H} c_e \le 2 \cdot \text{OPT}_{LP} \le 2 \cdot \text{OPT}
\]
which proves that Jain's iterative rounding is a 2-approximation.

\section{Applications}
\begin{itemize}
    \item \textbf{Survivable Network Design:} Laying communication links that can survive up to $k$ simultaneous link failures by maintaining $k$ disjoint paths between key nodes.
    \item \textbf{Power Grid Design:} Constructing high-capacity power lines that ensure alternative routes during generator outages.
    \item \textbf{Disaster Recovery Routing:} Deploying redundant emergency networks connecting key rescue centers.
\end{itemize}

\subsection*{Real Example: High-Availability Telecommunication Backbones}
\textbf{Scenario:} A telecom operator (e.g. AT\&T) connects regional routing hubs. Certain hubs (e.g., Chicago and New York) require a high-availability channel with 3 edge-disjoint paths to survive up to 2 simultaneous fiber cuts. Other hubs only require 1 path.
\textbf{Model:} Steiner Network. Hubs are vertices, connectivity requirements are $r_{uv} \in \{0, 1, 3\}$. Iterative rounding finds the optimal network.
\textbf{Why Iterative Rounding matters:} Dual-based schemas or randomized rounding cannot handle high disjointness constraints ($r_{uv} \ge 3$) efficiently. Jain's iterative rounding algorithm handles general integer demands directly, ensuring a cost-efficient layout with a strict mathematical guarantee of being within $2 \times$ of the global optimum.

\section{Tight Example: The Multi-Path Cycle}
Consider a 4-node graph with edges $(0, 1), (1, 2), (2, 3), (3, 0)$ of costs $2, 3, 2, 3$, and a diagonal edge $(0, 2)$ of cost 4. Let the requirements be $r(0, 2) = 2$ and $r(1, 3) = 1$.
The optimal integer solution is to select edges $(0, 1), (1, 2), (2, 3), (3, 0)$ with total cost $2 + 3 + 2 + 3 = 10$.
Our iterative rounding algorithm solves the LP, rounds the diagonal edge $(0, 2)$ (since its LP value is $\ge 0.5$), and then rounds $(0, 1), (1, 2), (2, 3)$, yielding a total cost of $2 + 3 + 2 + 4 = 11$.
The approximation ratio is $11/10 = 1.1$, which is well within the factor 2 performance bound.

\section{Exercises}
\begin{exercise}
Prove that if all connectivity requirements $r_{uv}$ are at most 1, the Steiner Network problem is equivalent to the Steiner Forest problem.
\end{exercise}
\begin{exercise}
Construct a basic feasible solution for a triangle graph with connectivity requirements 1, and show that at least one edge has $x_e \ge 1/2$.
\end{exercise}

